\documentclass[10pt,a4paper]{amsart}
\usepackage[margin=19mm]{geometry}
\usepackage{amsmath,amssymb,mathtools}
\usepackage[scr=rsfs,cal=euler]{mathalfa}
\usepackage{xfrac}
\usepackage[unicode,hidelinks,bookmarks=false]{hyperref}
\newcommand{\ie}{i.e.\ }
\newcommand{\cf}{cf.\ }
\newcommand{\resp}{respectively\ }
\newcommand{\Subsection}{\subsection}
\providecommand{\nobreakdash}{\nobreak\hbox{-}}
\setlength{\emergencystretch}{2em}
\multlinegap0pt
\usepackage{amssymb}
\usepackage{mathtools}
\usepackage{braket}
\usepackage[shortlabels]{enumitem}
\usepackage{xcolor}
\usepackage{csquotes}
\usepackage{dsfont}
\usepackage[normalem]{ulem}
\newtheorem{lemma}{Lemma}
\DeclareMathOperator{\Ad}{Ad}
\newcommand{\comm}{\mathrm{Comm}}
\title{Minimal model for confined operator dynamics in Floquet quantum circuits}
\author{Marcell D. Kov\'acs, Christopher J. Turner, Llu\'is Masanes}
\date{Source: arXiv:2602.06913v4 (CC BY 4.0)}
\begin{document}
\maketitle

\noindent\emph{Source and licence.} Minimal model for confined operator dynamics in Floquet quantum circuits. Version arXiv:2602.06913v4 (CC BY 4.0), distributed by arXiv under the Creative Commons Attribution 4.0 International licence, http://creativecommons.org/licenses/by/4.0/, as recorded in the arXiv licence metadata for this exact version and shown on the abstract page https://arxiv.org/abs/2602.06913v4. Full licence notice: \texttt{licenses/arxiv-2602.06913-CC-BY-4.0.txt}.\par\smallskip

\noindent\textit{Editorial scope.} Counted unit: the article's lemma lemma:U\_unique\_decomposition (source0.tex lines 1034--1037). The article's complete original proof of it is delivered with it (source0.tex lines 1038--1068, one \texttt{proof} environment of 31 lines). No mathematics is written, repaired or re-derived: the statement and the proof are the article's own lines, in the article's own notation, and no retained line is substituted or re-flowed.  No further source-local context is needed: every cross-reference of every retained line resolves inside the two delivered spans.  Nothing was omitted from any retained span, so the package carries no omission record.  No retained line contains a citation, so the wrapper carries no bibliography and no bibliography entry.  No retained line contains a figure environment, an image-inclusion command or drawing code, so no image is delivered and the package declares no asset dependency.  Editorial formatting only, in addition: an AMS \texttt{amsart} preamble that restates the article's own declarations and macros used by the retained lines, namely the environments \texttt{lemma} on separate counters, together with the standard packages those lines need.  The retained environments print in this document as Lemma 1 in order of appearance, whereas the article prints the same items with the numbering of its own full document.  The complete native source of the article as distributed in the arXiv e-print for this exact version is bundled unchanged as \texttt{sources/194139/source0.tex}.
\begin{lemma}[Representative elements of the coset space]  
    Within every coset in $\mathrm{U}(\mathcal{A})/\Sigma$ there is a unique representative element which leaves $Z(\mathcal{A})$ invariant, together these form a set of representatives $V$. This means that any automorphism $U \in \mathrm{U}(\mathcal{A})$ may be uniquely decomposed into a representative element $v \in V$ and a permutation $\sigma \in \Sigma$ such that $U = v \sigma$.
    \label{lemma:U_unique_decomposition}
\end{lemma}

\begin{proof}
  Let $U \in \mathrm{U}(\mathcal{A})$.
  Since $\mathrm{U}(\mathcal{A})$ is the unitary automorphism group of $\mathcal{A}$, the image of $Z(\mathcal{A})$ under $U$ is $Z(\mathcal{A})$ itself.
  Recall that the centre $Z(\mathcal{A})$ has a basis of projection operators, $\Pi_i$ for $i \in I$, corresponding to each term of the direct sum decomposition of $\mathcal{H}$.
  Examine the action of $U$ on this basis,
  \begin{align}
      \Ad_U : \Pi_i \mapsto \sum_{j \in J} \Pi_j & \text{ for some $J \subseteq I$.}
  \end{align}
  since it must preserve the eigenvalues of $\Pi_i$.
  If we take some element of the basis $\Pi_i$ such that the trace of $\Pi_i$ is minimal, then $J$ consist of a single element $j$.
  Since every other element of the basis is orthogonal to $\Pi_i$, the image of every other element is also orthogonal to the image of $\Pi_i$ so we can remove $i$ from the domain basis set and remove $j$ from the codomain basis set and proceed inductively.
  This constructs a permutation $\sigma$ acting on $I$ by mapping each $i$ to a corresponding element $j$.

  % \sigma \in \Sigma
  Now we must show that $\sigma$ is in $\Sigma$.
  The dimensions of $\mathcal{A}$ and $\comm(\mathcal{A})$ as algebras within an element $i$ of the decomposition are $D_i^2 = \dim (\Pi_i \mathcal{A} \Pi_i)$ and $E_i^2 = \dim(\Pi_i \comm(\mathcal{A}) \Pi_i)$.
  We can relate the dimensions within $i$ and $j$ as follows,
  \begin{align}
      \Pi_j \mathcal{A} \Pi_j
      &= \Pi_j U \mathcal{A} U^\dagger \Pi_j \\
      &= (U \Pi_i U^\dagger) U \mathcal{A} U^\dagger (U \Pi_i U^\dagger) \\
      &= U (\Pi_i \mathcal{A} \Pi_i) U^\dagger
  \end{align}
  which has the same dimension as $\Pi_i \mathcal{A} \Pi_i$ since $U$ is unitary.
  % Since $\Ad_{U}$ preserves the dimension of $\mathcal{A}$ and $\mathcal{A'}$ as algebras, we have that $D_i^2 = D_j^2$ and $E_i^2 = E_j^2$.
  Hence, $D_i^2 = D_j^2$ and $E_i^2 = E_j^2$.
  Therefore, the permutation is in the appropriate subgroup $\Sigma$ defined earlier.

  % Proof of Uniqueness
  Uniqueness of this representative element follows because the only element of $\Sigma$ which leaves $Z(\mathcal{A})$ fixed is the identity element.
\end{proof}

\end{document}
